% !TEX program = lualatex
\documentclass[aspectratio=43,professionalfonts,handout]{beamer}
\usepackage{beamer_template}

\newcommand{\LectureNum}{07}
\newcommand{\LectureTitle}{２次関数と２次不等式}
\newcommand{\TextPages}{pp.\,84-85}
\newcommand{\NextTopic}{問題演習}
\newcommand{\NextPages}{pp.\,86-87（練習問題1-A, 1-B）}
\newcommand{\CourseName}{基礎数学B}
\renewcommand{\TermName}{前期}

\begin{document}

% -----------------------------------------------
% スライド1: 表紙＋目標＋用語・定義
% -----------------------------------------------
\begin{frame}{\CourseName\quad 第\LectureNum 回「\LectureTitle 」}
  {\small \TermName\quad \TeacherName\quad ／\quad 教科書 \TextPages}

  \begin{block}{用語・定義}
  \begin{description}[labelwidth=5.5em]
    \item[\kw{2次不等式}] $ax^2+bx+c>0$ や $ax^2+bx+c<0$ の形の不等式
    \item[\kw{解き方}] 放物線と $x$ 軸の位置関係から $y>0$（$y<0$）となる $x$ の範囲を読み取る
  \end{description}
  \end{block}
\end{frame}

% -----------------------------------------------
% スライド2: 公式・性質
% -----------------------------------------------
\begin{frame}{公式・性質}
  \begin{block}{}
  \begin{itemize}
    \item $D>0$（解 $\alpha<\beta$）：$ax^2+bx+c>0$ の解は $x<\alpha$ または $x>\beta$
    \item $D=0$：$ax^2+bx+c>0$ の解は $x\neq\alpha$（$\alpha$ 以外のすべての実数）
    \item $D<0$：$ax^2+bx+c>0$ の解はすべての実数
    \item $<0$, $\leq 0$, $\geq 0$ は上記の符号・等号を変えて対応
  \end{itemize}
  \end{block}
\end{frame}

% -----------------------------------------------
% スライド3: 例1→問1
% -----------------------------------------------
\begin{frame}{例1→問1}
  \begin{exampleblock}{例題7) p.84（3パターン各1問）}
    （板書で解説）
  \end{exampleblock}

  \vfill

  \begin{block}{問13) p.85（4問）\quad 自分で解いてみよう}
    （ノートに解くこと）
  \end{block}
\end{frame}

% -----------------------------------------------
% スライド4: 追加演習
% -----------------------------------------------
\begin{frame}{追加演習}
  \begin{block}{練習1-A 問8・問9) p.86\quad 自分で解いてみよう}
    （ノートに解くこと）
  \end{block}
\end{frame}

% -----------------------------------------------
% まとめ
% -----------------------------------------------
\begin{frame}{まとめ}
  \begin{itemize}
    \item 2次不等式は放物線と $x$ 軸の位置関係をグラフで確認して解く
    \item $D>0$・$D=0$・$D<0$ の3パターンで解の形が異なる
    \item 「解なし」「すべての実数」が答えになるケースを見落とさない
  \end{itemize}

  \vfill
  {\small 基本問題：155, 164}
\end{frame}

\end{document}
